A key feature of projective representations is that we also have the following canonical
decomposition whose proof can be obtained in a similar way as in the case of regular
representations.

\begin{Th}[canonical decomposition]\label{can,decom}
Suppose that $\alpha$ is a
normalized $2$-cocycle of a finite group $G$ with values in $\S^{1}$. Let $W$ be a
finite-dimensional $\alpha$-representation. Then the assignment
\begin{align*}
\gamma\colon\ \bigoplus_{[\rho]\in \Irr^{\alpha}(G)}V_{\rho}\otimes \Hom_{G}(V_{\rho},W)
&\rightarrow W,
\\
v\otimes f&\mapsto f(v)
\end{align*}
defines an isomorphism of $\alpha$-representations.
\end{Th}

Suppose now that $\alpha$ is a normalized 2-cocycle on $G$ with values in $\S^{1}$.
We can associate to $\alpha$ a central extension of $G$ by $\S^{1}$ in the following way.
As a set define
\begin{gather*}%\label{defcentralext}
\widetilde{G}_{\alpha}=\big\{(g,z)\mid g\in G, z\in \S^{1}\big\}.
\end{gather*}
The product structure in $\widetilde{G}_{\alpha}$ is given by
the assignment
\begin{gather*}
(g_{1},z_{1})(g_{2},z_{2}):=\left(g_{1}g_{2},\alpha(g_{1},g_{2})z_{1}z_{2}\right).
\end{gather*}
This way $\widetilde{G}_{\alpha}$ is a compact Lie group that fits
into a central extension
\begin{gather*}
1\rightarrow \S^{1}\rightarrow
\widetilde{G}_{\alpha}\rightarrow G\rightarrow 1.
\end{gather*}
Let $\rho\colon G\to {\rm GL}(V)$ be an $\alpha$-representation of $G$. If we define
$\tilde{\rho}\colon \widetilde{G}_{\alpha}\to {\rm GL}(V)$
by $\tilde{\rho}(g,z)=z\rho(g)$ then $\tilde{\rho}$ defines a representation of
$\widetilde{G}_{\alpha}$ on which the central factor $\S^{1}$ acts by multiplication of~scalars.
Conversely, if $\tilde{\rho}\colon \widetilde{G}_{\alpha}\to {\rm GL}(V)$ is a representation of
$\widetilde{G}_{\alpha}$ on which the central factor~$\S^{1}$ acts by multiplication of scalars
then the function $\rho\colon G\to {\rm GL}(V)$ given by $\rho(g)=\tilde{\rho}(g,1)$ defines an
$\alpha$-representation of $G$. The above assignment defines a one to one correspondence between
$\alpha$-representations of $G$ and representations of $\widetilde{G}_{\alpha}$ on which the
central factor $\S^{1}$ acts by multiplication of scalars. Via this correspondence we will
switch back and forth between $\alpha$-representations of $G$ and representations of
$\widetilde{G}_{\alpha}$ on which the central factor $\S^{1}$ acts by multiplication of scalars
without explicitly mentioning it. In addition, if $V_{1}$ and $V_{2}$ are two
$\alpha$-repre\-sentations of $G$ then having a map of projective representations
$f\colon V_{1}\to V_{2}$ is equivalent to~having a linear map
$f\colon V_{1}\to V_{2}$ that is $\widetilde{G}_{\alpha}$-equivariant. With this correspondence
in mind we~will also identify $\Hom_{G}(V_{1},V_{2})$ with
$\Hom_{\widetilde{G}_{\alpha}}(V_{1},V_{2})$ without explicitly mentioning it.

\subsection{Cocycles and projective representations}\label{section cocycles}

Suppose now that $A$ is a normal subgroup of a finite group $G$ so that we have a short exact
sequence
\begin{gather*}
1\rightarrow A\rightarrow G\stackrel{\pi}{\rightarrow} Q\rightarrow 1.
\end{gather*}
Assume that $\alpha$ is a normalized $2$-cocycle of $G$ with values in $\S^{1}$,
by restriction we can also see $\alpha$ as a cocycle defined on $A$.
Let us define an action of $G$ on the set $\Irr^{\alpha}(A)$ in the following way.
Given $\rho\colon A\rightarrow U(V_{\rho})$ an $\alpha$-representation and $g\in G$ we define
$g\cdot \rho\colon A\rightarrow U(V_{\rho})$ so that if~$a\in A$ we have:
\begin{gather}\label{def:act}
g\cdot \rho(a)=\alpha(g^{-1}a,g)\alpha(g,g^{-1}a)^{-1}\rho(g^{-1}ag)\in U(V_{\rho}).
\end{gather}

\begin{Pro}
The above assignment defines a left action of $G$ on $\Irr^{\alpha}(A)$. Furthermore, for all
$b\in A$, we have that $b\cdot\rho\cong\rho$ so that the action of $G$ on $\Irr^{\alpha}(A)$
factors to an action of~$Q=G/A$ on $\Irr^{\alpha}(A)$.
\end{Pro}

\begin{proof}
First we show that $g\cdot \rho$ is an $\alpha$-representation
of $A$. Indeed, for all $g\in G$ and $a,b \in A$ we have
\begin{align*}
(g\cdot \rho(a))(g\cdot\rho(b))&=
\big(\alpha\big(g^{-1}a,g\big)\alpha\big(g,g^{-1}a\big)^{-1}\rho\big(g^{-1}ag\big)\big)
\big(\alpha\big(g^{-1}b,g\big)\alpha\big(g,g^{-1}b\big)^{-1}\rho\big(g^{-1}bg\big)\big)
\\
&=\alpha\big(g^{-1}a,g\big)\alpha\big(g,g^{-1}a\big)^{-1}\alpha\big(g^{-1}b,g\big)\alpha\big(g,g^{-1}b\big)^{-1}
\alpha\big(g^{-1}ag,g^{-1}bg\big)
\\
&\phantom{=}\,\times\rho\big(g^{-1}abg\big)
=\alpha\big(g,g^{-1}a\big)^{-1}\alpha\big(g^{-1}ab,g\big)\alpha\big(g^{-1}a,b\big)\rho\big(g^{-1}abg\big)
\\
&=\alpha(a,b)\alpha\big(g^{-1}ab,g\big)\alpha\big(g,g^{-1}ab\big)^{-1}\rho\big(g^{-1}abg\big)\\
&=\alpha(a,b) (g\cdot \rho(ab) ).
\end{align*}
The above equalities are obtained using the 2-cocycle equation for $\alpha$.
Now, we show that this definition satisfies the axioms of an action. If
$\rho\colon A\rightarrow U(V_{\rho})$ is an $\alpha$-representation,
as $\alpha$ is a~normalized cocycle,
\begin{gather*}
1\cdot \rho(a)=\alpha(a,1)\alpha(1,a)^{-1}\rho(a)=\rho(a).
\end{gather*}
Moreover, given $g,h\in G$ and $a\in A$, we have
\begin{align*}
g\cdot(h\cdot \rho)(a)&=\alpha\big(g^{-1}a,g\big)\alpha\big(g,g^{-1}a\big)^{-1}(h\cdot
\rho)\big(g^{-1}ag\big)
\\
&=\alpha\big(g^{-1}a,g\big)\alpha\big(g,g^{-1}a\big)^{-1}\alpha\big(h^{-1}g^{-1}ag,h\big)
\alpha\big(h,h^{-1}g^{-1}ag\big)^{-1}\rho\big(h^{-1}g^{-1}agh\big)
\end{align*}
and
\begin{align*}
(gh)\cdot \rho(a)&=\alpha\big(h^{-1}g^{-1}a,gh\big)\alpha\big(gh,h^{-1}g^{-1}a\big)^{-1}\rho\big((gh)^{-1}a(gh)\big)
\\
&=\alpha\big(h^{-1}g^{-1}a,gh\big)\alpha\big(gh,h^{-1}g^{-1}a\big)^{-1}\rho\big(h^{-1}g^{-1}agh\big).
\end{align*}
Manipulating the cocycle equation for $\alpha$ it can be proved that
\begin{align*}
\alpha\big(g^{-1}a,g\big)&\alpha\big(g,g^{-1}a\big)^{-1}\alpha\big(h^{-1}g^{-1}ag,h\big)\alpha\big(h,h^{-1}g^{-1}ag\big)^{-1}
\\
&=\alpha\big(h^{-1}g^{-1}a,gh\big)\alpha\big(gh,h^{-1}g^{-1}a\big)^{-1}.
\end{align*}
This implies that for all $a\in A$
\begin{gather*}
g\cdot(h\cdot \rho)(a)=(gh)\cdot \rho(a).
\end{gather*}
Finally, for $a, b\in A$ expanding and using the cocycle equation we obtain
\begin{gather*}
b\cdot\rho(a)=\rho(b)^{-1}\rho(a)\rho(b)
\end{gather*}
and thus $b\cdot\rho\cong\rho$ as $\alpha$-representations.
\end{proof}

\begin{Rem}The action of $G$ on the set $\Irr^{\alpha}(A)$ given by equation~(\ref{def:act}) can be
described in an alternative way as follows. Suppose that $\alpha$ is a normalized $2$-cocycle of
$G$ with values in~$\S^{1}$ and let $\widetilde{G}_{\alpha}$ be the central extension corresponding
to the cocycle $\alpha$. If~$\rho\colon A\rightarrow U(V_{\rho})$ is an~$\alpha$-representation of $A$
then as explained before we have an associated representation \mbox{$\tilde{\rho}\colon \tilde{A}_{\alpha}\to U(V_{\rho})$}.
Given $(g,z)\in \widetilde{G}_{\alpha}$ we obtain the $\tilde{A}_{\alpha}$-representation
$(g,z)\cdot \tilde{\rho}$ defined by
\begin{gather*}
(g,z)\cdot \tilde{\rho}(a,w)=\tilde{\rho}\big((g,z)^{-1}(a,w)(g,z)\big).
\end{gather*}
This is a well-defined $\tilde{A}_{\alpha}$-representation such that the central factor
$\S^{1}$ acts by multiplication of scalars. Interpreting this $\tilde{A}_{\alpha}$-representation as an
$\alpha$-representation and using the cocycle identities we obtain equation~(\ref{def:act}). Explicitly,
for $g\in G$ and $a\in A$ we have
\begin{align*}
g\cdot \rho(a):\!&=(g,1)\cdot \tilde{\rho}(a,1)= \tilde{\rho}\big(g^{-1}ag,\alpha\big(g^{-1}a,g\big)\alpha\big(g,g^{-1}a\big)^{-1}\big)
\\
&=\alpha\big(g^{-1}a,g\big)\alpha\big(g,g^{-1}a\big)^{-1} \rho\big(g^{-1}ag\big).
\end{align*}
\end{Rem}

%\medskip

As above assume that $A$ is a normal subgroup of a group $G$ and $Q=G/A$. Fix an assignment
$\sigma \colon Q \rightarrow G$ such that $\pi(\sigma(q)) = q$ for all $q \in Q$ with $\sigma(1)
= 1$. We remark that the map $\sigma$ is only a set theoretical map so in particular it does
not necessarily have to be a group homomorphism.

Suppose that $\rho\colon A \rightarrow U(V_{\rho})$ is a complex irreducible $\alpha$-representation
with the property that $g\cdot \rho$ is isomorphic to $\rho$ for every $g\in G$
(under the action defined in equation~(\ref{def:act})). Under this assumption,
as $\sigma(q)\cdot \rho\cong\rho$
we can find an element $M_{q}\in
U(V_{\rho})$ for each $q \in Q$ such that
\begin{gather*}%\label{ecu:Mq}
	\sigma(q)\cdot \rho(a)=M_{q}^{-1}\rho(a)M_{q}.
\end{gather*}
This means that
\begin{gather*}
\alpha\big(\sigma(q)^{-1}a,\sigma(q)\big)\alpha\big(\sigma(q),\sigma(q)^{-1}a\big)^{-1}
\rho\big(\sigma(q)^{-1}a\sigma(q)\big)=M_{q}^{-1}\rho(a)M_{q}
\end{gather*}
for all $a\in A$. We can choose $M_{1} = 1$ as $\sigma(1)=1$ and $\sigma(1)\cdot\rho=\rho$.
Let $\widetilde{A}_{\alpha}$ be the central extension of $A$ by $\S^{1}$
associated to the cocycle $\alpha$ and $\tilde{\rho}\colon \widetilde{A}_{\alpha}\to U(V_{\rho})$ the
corresponding representation. Remember that $\tilde{\rho}(a,z)=z\rho(a)$ so we can define
\begin{gather*}
\sigma(q)\cdot \tilde{\rho}(a,z):=z (\sigma(q)\cdot \rho(a)).
\end{gather*}
Therefore
\begin{gather*}%\label{ecu:Mrhot}
\sigma(q)\cdot \tilde{\rho}(a,z)=z\big(M_{q}^{-1}\rho(a)M_{q}\big)
=M_{q}^{-1}z\rho(a)M_{q}=M_{q}^{-1}\tilde{\rho}(a,z)M_{q}.
\end{gather*}
Define $\chi\colon Q\times Q \rightarrow A$ by the equation
\begin{gather*}
\chi(q_{1},q_{2})=\sigma(q_{1}q_{2})^{-1}\sigma(q_{1})\sigma(q_{2})%\label{ecu:chi}.
\end{gather*}
Note that $\chi(q_{1},q_{2})$ belongs to $A$ since
$\pi(\chi(q_{1},q_{2}))=1$ and the map $\chi$ is normalized in
the sense that $\chi(q_{1},q_{2})=1$ whenever $q_{1}=1$ or
$q_{2}=1$. In addition, define $\tau\colon Q\times Q\rightarrow \S^{1}$ by
\begin{gather*}
\tau(q_{1},q_{2})=\alpha\big(\sigma(q_{1}q_{2})^{-1},\sigma(q_{1})\sigma(q_{2})\big)
\alpha\big(\sigma(q_{1}q_{2}),\sigma(q_{1}q_{2})^{-1}\big)^{-1}\alpha(\sigma(q_{1}),\sigma(q_{2}))\nonumber
\\ \hphantom{\tau(q_{1},q_{2})}
{}=\alpha\big(\sigma(q_{1}q_{2}),\chi(q_{1},q_{2})\big)^{-1}\alpha(\sigma(q_{1}),\sigma(q_{2})).
%\label{ecu:tau}
\end{gather*}
Now, for $q_{1}$, $q_{2}$ in $Q$ we notice that the element
\begin{gather*}
\tilde{\rho}(\chi(q_{1},q_{2}),\tau(q_{1},q_{2}))M_{q_{2}}^{-1}M_{q_{1}}^{-1}M_{q_{1}q_{2}}
=\tau(q_{1},q_{2})\rho(\chi(q_{1},q_{2}))M_{q_{2}}^{-1}M_{q_{1}}^{-1}
M_{q_{1}q_{2}}
\end{gather*}
belongs to the center $Z(U(V_{\rho}))\cong \S^{1}$. Define the map
$\beta_{\rho,\alpha}\colon Q\times Q \rightarrow
\S^{1}$ by the equation
\begin{gather}
\label{defQbeta}
\beta_{\rho,\alpha}(q_{1},q_{2}):
=\tilde{\rho}(\chi(q_{1},q_{2}),\tau(q_{1},q_{2}))M_{q_{2}}^{-1}M_{q_{1}}^{-1}M_{q_{1}q_{2}}.
\end{gather}

\begin{Lem}\label{twistedcocycle}
If $\rho\colon A \rightarrow U(V_{\rho})$ is a complex
irreducible $\alpha$-representation such that $g\cdot \rho\cong \rho$ for every
$g\in G$ then the map $\beta_{\rho,\alpha}\colon Q\times Q\rightarrow \S^{1}$ defines a
normalized $2$-cocycle on $Q$ with values in~$\S^{1}$.
\end{Lem}

\begin{proof}
If either $q_{1}=1$ or $q_{2}=1$ we have that $\tau(q_{1},q_{2})=1$ as $\alpha$ is normalized.
Also, as $\chi$ is normalized we have $\chi(q_{1},q_{2})=1$. Since we are choosing
$\sigma(1)=1$ and $M_{1}=1$ it follows that $\beta_{\rho,\alpha}(q_{1},q_{2})=1$ if either
$q_{1}=1$ or $q_{2}=1$. Therefore either $\beta_{\rho,\alpha}$ is normalized. To finish we need
to prove that for every $q_{1},q_{2},q_{3}\in Q$ we have
\begin{gather*}
\beta_{\rho,\alpha}(q_{1},q_{2}q_{3})\beta_{\rho,\alpha}(q_{2},q_{3})
=\beta_{\rho,\alpha}(q_{1}q_{2},q_{3})\beta_{\rho,\alpha}(q_{1},q_{2}).
\end{gather*}
To see this note that as $\beta_{\rho,\alpha}(q_{1},q_{2})$ belongs to $\S^{1}$ we
have that
\begin{gather*}
M_{q_{1}q_{2}}
=\beta_{\rho,\alpha}(q_{1},q_{2})M_{q_{1}}M_{q_{2}}
\tilde{\rho}(\chi(q_{1},q_{2}),\tau(q_{1},q_{2}))^{-1}.
\end{gather*}
Therefore, for $q_1, q_2$ and $q_3$ in $Q$ we have,
\begin{align*}
M_{q_{1}q_{2}q_{3}}&=M_{q_{1}(q_{2}q_{3})}
=\beta_{\rho,\alpha}(q_{1},q_{2}q_{3})M_{q_{1}}
M_{q_{2}q_{3}}\tilde{\rho}(\chi(q_{1},q_{2}q_{3}),\tau(q_{1},q_{2}q_{3}))^{-1}
\\
&=\beta_{\rho,\alpha}(q_{1},q_{2}q_{3})M_{q_{1}}
\beta_{\rho,\alpha}(q_{2},q_{3})M_{q_{2}}M_{q_{3}}
\tilde{\rho}(\chi(q_{2},q_{3}),\tau(q_{2},q_{3}))^{-1}
\\
&\phantom{=}\times\tilde{\rho}(\chi(q_{1},q_{2}q_{3}),
\tau(q_{1},q_{2}q_{3}))^{-1}
=\beta_{\rho,\alpha}(q_{1},q_{2}q_{3})\beta_{\rho,\alpha}(q_{2},q_{3})M_{q_{1}}M_{q_{2}}M_{q_{3}}
\\
&\phantom{=}\times\big[\tilde{\rho}(\chi(q_{1},q_{2}q_{3}),\tau(q_{1},q_{2}q_{3}))
\tilde{\rho}(\chi(q_{2},q_{3}),\tau(q_{2},q_{3}))\big]^{-1}.
\end{align*}
In a similar way, writing $M_{q_{1}q_{2}q_{3}}=M_{(q_{1}q_{2})q_{3}}$ and expanding we obtain
\begin{align*}
M_{q_{1}q_{2}q_{3}}&=\beta_{\rho,\alpha}(q_{1}q_{2},q_{3})
\beta_{\rho,\alpha}(q_{1},q_{2})M_{q_{1}}M_{q_{2}}M_{q_{3}}
\\
&\phantom{=}\times\big[\tilde{\rho}(\chi(q_{1}q_{2},q_{3}),\tau(q_{1}q_{2},q_{3}))
\tilde{\rho}(\chi(q_{1},q_{2}),\tau(q_{1},q_{2}))\big]^{-1}.
\end{align*}	
This implies that
\begin{gather*}
\beta_{\rho,\alpha}(q_{1},q_{2}q_{3})\beta_{\rho,\alpha}(q_{2},q_{3})
=\beta_{\rho,\alpha}(q_{1}q_{2},q_{3})\beta_{\rho,\alpha}(q_{1},q_{2})
\end{gather*}
as we wanted to prove.
\end{proof}

%
%*************************************** SECTION DECOMPOSITION ***********************************
%

\section[Decomposition of twisted equivariant K-theory]
{Decomposition of twisted equivariant $\boldsymbol K$-theory}\label{Section 3}

In this section we use the canonical decomposition of vector bundles to show that,
under some hypothesis, the $\alpha$-twisted equivariant $K$-theory $^{\alpha}K^{\ast}_{G}(X)$ of a
$G$-space $X$ can be decomposed as a direct sum of twisted equivariant K-theories parametrized
by the orbits of the action of $G$ on~$\Irr^{\alpha}(A)$ constructed on the previous section.

We start by recalling the definition of $\alpha$-twisted equivariant $K$-theory that we will use.
We~fol\-low the treatment used in~\cite[Section 7.2]{AdemRuan}.
Assume that $G$ is a finite group acting on a compact and Hausdorff space $X$.
Let $\alpha$ be a normalized 2-cocycle on $G$ with values in $\S^{1}$. Consider
\begin{gather*}
1\rightarrow \S^{1}\rightarrow \widetilde{G}_{\alpha}\rightarrow G\rightarrow 1
\end{gather*}
the corresponding central extension. Observe that the action of $G$ on $X$ can be
extended to an~action of $\widetilde{G}_{\alpha}$ in such a way that the central factor $\S^{1}$
acts trivially.

\begin{Def}\label{usualdeftwisted}
The $\alpha$-twisted $G$-equivariant $K$-theory of $X$, denoted by $^{\alpha}K_{G}^{0}(X)$,
is defined as the Grothendieck group of the set of isomorphism classes of
$\widetilde{G}_{\alpha}$-equivariant vector bundles over $X$ on which
$\S^{1}$ acts by multiplication of scalars on the fibers. For $n>0$ the twisted groups
$^{\alpha}K_{G}^{n}(X)$ are defined as
$^{\alpha}\tilde{K}_{G}^{0}(\Sigma^{n}X_{+})$, where as usual $X_{+}$
denotes the space $X$ with an added base point.
\end{Def}

To obtain the desired decomposition of twisted equivariant $K$-groups we are going to construct an equivalent
formulation for such twisted K-groups following the work in~\cite{GomezUribe}. For this suppose
that we have a short exact sequence of finite groups
\begin{gather*}
1\rightarrow A\rightarrow G\stackrel{\pi}{\rightarrow} Q\rightarrow 1.
\end{gather*}
Assume that $G$ acts on a compact and Hausdorff space $X$ in such a way that $A$ acts trivially.
Fix $\alpha\in Z^{2}\big(G,\S^{1}\big)$ a normalized 2-cocycle. Notice that by restriction we can
see any $\widetilde{G}_{\alpha}$-equivariant vector bundle over $X$ as an
$\widetilde{A}_{\alpha}$-equivariant vector bundle, where $\widetilde{A}_{\alpha}$ denotes the
central extension associated to the cocycle $\alpha$ seen as a cocycle defined on $A$.
Also, recall that associated to an~$\alpha$-representation $\rho\colon A \rightarrow U(V_{\rho})$ we have a representation
$\tilde{\rho}\colon \widetilde{A}_{\alpha} \rightarrow U(V_{\rho})$.

{\sloppy\begin{Def}
Suppose that $\rho\colon A \rightarrow U(V_{\rho})$ is a complex irreducible $\alpha$-representation.
A~$(\widetilde{G}_{\alpha}, \rho)$-equivariant vector bundle over $X$ is a
$\widetilde{G}_{\alpha}$-vector bundle on which the central factor $\S^{1}$ acts
by multiplication of scalars on $E$ and the map
\begin{align*}
\gamma\colon\ \V_{\rho}\otimes \Hom_{\widetilde{A}_{\alpha}}(\V_{\rho},E)&\rightarrow E,
\\
v\otimes f &\mapsto f(v)
\end{align*}
is an isomorphism of $\widetilde{A}_{\alpha}$-vector bundles on which $\S^{1}$ acts
by multiplication of scalars.
\end{Def}

}

In the above definition, if $\rho$ is an $\alpha$-representation of $A$ then $\V_{\rho}$
denotes the trivial $\widetilde{A}_{\alpha}$-vector bundle $\pi_{1}\colon X\times V_{\rho}\to X$.
Observe that if $p\colon E\to X$ is a $\widetilde{G}_{\alpha}$-vector bundle on which the central
factor $\S^{1}$ acts by multiplication then, as we are assuming that $A$ acts trivially
on $X$, it follows that for every $x \in X$ the fiber $E_{x}$ can be seen as a representation of
$\widetilde{A}_{\alpha}$ on which the central factor acts by scalar multiplication. Thus for every
$x \in X$ the fiber $E_{x}$ can be seen as an $\alpha$-representation of $A$.
With this point of view, a $\big(\widetilde{G}_{\alpha}, \rho\big)$-equivariant vector bundle is a
$\widetilde{G}_\alpha$-equivariant vector bundle $p\colon E\rightarrow X$ such that for every
$x \in X$ the fiber $E_{x}$ is an~$\alpha$-representation of $A$ isomorphic to a direct sum of
the $\alpha$-representation $\rho$. Let $\text{Vect}_{\widetilde{G}_{\alpha}, \rho}(X)$
denote the set of isomorphism classes of $\big(\widetilde{G}_{\alpha}, \rho\big)$-equivariant vector
bundles, where two $\big(\widetilde{G}_{\alpha}, \rho\big)$-equivariant vector bundles are isomorphic
if they are isomorphic as $\widetilde{G}_{\alpha}$-vector bundles. Notice that if
$E_{1}$ and $E_{2}$ are two $\big(\widetilde{G}_{\alpha}, \rho\big)$-equivariant vector bundles
then so is $E_{1}\oplus E_{2}$. Therefore $\text{Vect}_{\widetilde{G}_{\alpha}, \rho}(X)$
is a semigroup. Following~\cite[Definition~2.2]{GomezUribe} we have the next definition.

\begin{Def}\label{def,grhoalfa}
Assume that $G$ acts on a compact space $X$ in such a way that $A$ acts trivially on $X$
and let $\alpha$ be a normalized 2-cocycle $\alpha\in Z^{2}\big(G,\S^{1}\big)$. We define
$K^{0}_{\widetilde{G}_{\alpha}, \rho}(X)$, the $\big(\widetilde{G}_{\alpha}, \rho\big)$-equivariant
$K$-theory of $X$, as the Grothendieck construction applied to
$\text{Vect}_{\widetilde{G}_{\alpha}, \rho}(X)$. For $n>0$ the group
$K^{n}_{\widetilde{G}_{\alpha}, \rho}(X)$ is defined as
$\tilde{K}^{0}_{\widetilde{G}_{\alpha}, \rho}(\Sigma^{n}X_{+})$.
\end{Def}

As our next step we show that the previous definition can be described using the usual definition
of twisted equivariant $K$-theory provided in Definition~\ref{usualdeftwisted}.
For this suppose that $\alpha$ is a normalized 2-cocycle of $G$ with values in $\S^{1}$.
As above assume that $A$ is a normal subgroup of $G$ and let $Q=G/A$. Let $\rho$ be an
irreducible $\alpha$-representation such that $g\cdot\rho\cong\rho$ for all $g \in G$. Fix
a set theoretical section $\sigma\colon Q\to G$ such that $\sigma(1)=1$ as in the previous section.
We can extend $\sigma$ to obtain a map $\tilde{\sigma}\colon Q\to \widetilde{G}_{\alpha}$ by defining
$\tilde{\sigma}(q)=(\sigma(q),1)\in \widetilde{G}_{\alpha}$. Let $\beta_{\rho,\alpha}$ be the
$2$-cocycle defined on $Q$ with values in $\S^{1}$ constructed
in equation~(\ref{defQbeta}). With this cocycle we can consider the central extension
\begin{gather*}
1\rightarrow \S^{1}\rightarrow \widetilde{Q}_{\beta_{\rho,\alpha}}
\rightarrow Q\rightarrow 1.
\end{gather*}
With this in mind we have the following generalization of~\cite[Theorem 2.1]{GomezUribe}.

\begin{Th}\label{teo,bij}
Suppose that $\rho$ is an irreducible $\alpha$-representation such that $g\cdot\rho\cong\rho$
for all $g \in G$. Let $X$ be a $G$-space such that $A$ acts trivially on $X$.
If $p\colon E\rightarrow X$ is a $\big(\widetilde{G}_{\alpha}, \rho\big)$-equivariant vector bundle, then
$\Hom_{\widetilde{A}_{\alpha}}(\V_{\rho},E)$ has the
structure of a $\widetilde{Q}_{\beta_{\rho,\alpha}}$-vector bundle on which
the central factor~$\S^{1}$ acts by multiplication of
scalars. Moreover, the assignment
\begin{gather*}
[E]\rightarrow
[\Hom_{\widetilde{A}_{\alpha}}(\V_{\rho},E)]
\end{gather*}
is a natural one to one correspondence between isomorphism
classes of $\big(\widetilde{G}_{\alpha},\rho\big)$-equivariant
vector bundles over $X$ and isomorphism classes of
$\widetilde{Q}_{\beta_{\rho,\alpha}}$-equivariant vector bundles over $X$
for which the central factor $\S^{1}$ acts by multiplication of scalars.
\end{Th}
\begin{proof} We are only going to provide a sketch of the proof as it follows the same
steps used in the proof of~\cite[Theorem 2.1]{GomezUribe}.

Suppose that $p\colon E \rightarrow X$ is a $\widetilde{G}_{\alpha}$-vector bundle. We give
$\Hom_{\widetilde{A}_{\alpha}}(\V_{\rho},E)$ an action of
	$\widetilde{Q}_{\beta_{\rho,\alpha}}$. Given $f \in
\Hom_{\widetilde{A}_{\alpha}}(\V_{\rho},E)_{x}$ and
$q\in Q$ we define $q \bullet f \in
\Hom_{\widetilde{A}_{\alpha}}(\V_{\rho},E)_{q\cdot x}$ by
\begin{gather*}
(q\bullet f)(v)=\tilde{\sigma}(q)\cdot f\big(M_{q}^{-1}v\big)
=(\sigma(q),1) \cdot f\big(M_{q}^{-1}v\big),
\end{gather*}
where $M_{q} \in U(V_{\rho})$ is the element chosen in the Section~\ref{section cocycles}.
It is easy to see that $q \bullet f$ is $\widetilde{A}_{\alpha}$-equivariant. Also,
if $q_{1}, q_{2}\in Q$ then $q_{1}\bullet
(q_{2} \bullet f)$ is such that for $v \in V_{\rho}$
\begin{align*}
q_{1}\bullet(q_{2}\bullet f)(v)&=\tilde{\sigma}(q_{1})
(q_{2}\bullet f)\big(M_{q_{1}}^{-1}v\big)=\tilde{\sigma}(q_{1})\tilde{\sigma}(q_{2})
f\big(M_{q_{2}}^{-1}M_{q_{1}}^{-1}v\big)\\
&=(\sigma(q_{1}q_{2}),1)(\chi(q_{1},q_{2}),\tau(q_{1},q_{2}))
f\big(M_{q_{2}}^{-1}M_{q_{1}}^{-1}v\big)\\
&=\tilde{\sigma}(q_{1}q_{2})f(\tilde{\rho}\big((\chi(q_{1},q_{2}),\tau(q_{1},q_{2}))
M_{q_{2}}^{-1}M_{q_{1}}^{-1}v\big)
\\
&=\tilde{\sigma}(q_{1}q_{2})
f\big(\beta_{\rho,\alpha}(q_{1},q_{2})M_{q_{1}q_{2}}^{-1}v\big)=\beta_{\rho,\alpha}(q_{1},q_{2}) \tilde{\sigma}(q_{1}q_{2})
f\big(M_{q_{1}q_{2}}^{-1}v\big)\\
&=\beta_{\rho,\alpha}(q_{1},q_{2})(q_{1}q_{2}\bullet f(v)).
\end{align*}
We conclude that
\begin{gather*}
q_{1}\bullet(q_{2}\bullet
f)(v)=\beta_{\rho,\alpha}(q_{1},q_{2})(q_{1}q_{2}\bullet f(v)).
\end{gather*}
The last equation allows us to define an action of
$\widetilde{Q}_{\beta_{\rho,\alpha}}$ on
$\Hom_{\widetilde{A}_{\alpha}}(\V_{\rho},E)$ as follows. If
$(q, z) \in Q_{\beta_{\rho,\alpha}}$ and $f \in
\Hom_{\widetilde{A}_{\alpha}}(\V_{\rho},E)_{x}$ define
\begin{gather*}
(q,z)\cdot f:=z(q\bullet f).
\end{gather*}
Thus if $v\in V_{\rho}$ then
\begin{gather*}
((q,z)\cdot f)(v)=z\big(\tilde{\sigma}(q)\cdot f\big(M_{q}^{-1}v\big)\big)=
\tilde{\sigma}(q)\cdot \big(zf\big(M_{q}^{-1}v\big)\big).
\end{gather*}
Unraveling the definitions, it is easy to see that this way
$\Hom_{\widetilde{A}_{\alpha}}(\V_{\rho},E)$ has the structure of
a~$\widetilde{Q}_{\beta_{\rho,\alpha}}$-equivariant vector bundle such that
the central factor $\S^{1}$ acts by multiplication of scalars.

Suppose now that $p \colon F \rightarrow X$ is a
$\widetilde{Q}_{\beta_{\rho,\alpha}}$-equivariant vector bundle over $X$
for which the central~$\S^{1}$ acts by multiplication of
scalars. Given $(g,z)\in \widetilde{G}_{\alpha}$, $f\in F_{x}$ and $v\in
V_{\rho}$ define
\begin{gather}
(g,z)\cdot (v\otimes f):=M_{\pi(g)}\tilde{\rho}\big(\tilde{\sigma}(\pi(g))^{-1}(g,z)\big)v\otimes
((\pi(g),1)\cdot f)\in (\V_{\rho}\otimes F)_{\pi(g)\cdot x}.
\label{ecu:ac2}
\end{gather}
The above assignment defines an action of $\widetilde{G}_{\alpha}$
on $\V_{\rho}\otimes F$ and $\V_{\rho}\otimes F$
becomes a $\widetilde{G}_{\alpha}$-vector bundle.
Moreover for $(a,z)\in \widetilde{A}_{\alpha}$, as $M_{1}=1$, we have
\begin{gather*}
(a,z)\cdot (v\otimes
f)=M_{1}\tilde{\rho}\big(\tilde{\sigma}(1)^{-1}(a,z)\big)v\otimes
(1,1)\cdot f=\tilde{\rho}(a,z)v\otimes f
\end{gather*}
so that $\widetilde{A}_{\alpha}$ acts on $\V_{\rho}\otimes
F$ by the representation $\tilde{\rho}$; that is,
$p\colon \V_{\rho}\otimes F\rightarrow X$ is a
$\big(\widetilde{G}_{\alpha},\rho\big)$-equivariant vector bundle
over~$X$.
	
Finally, assume that $p\colon E \rightarrow X$ is a
$\big(\widetilde{G}_{\alpha},\rho\big)$-equivariant vector bundle over $X$.
By definition the map
\begin{align*}
\gamma\colon \ \V_{\rho}\otimes
\Hom_{\widetilde{A}_{\alpha}}(\V_{\rho},E)&\rightarrow E,
\\
(v,f)&\mapsto f(v)
\end{align*}
is an isomorphism of $\widetilde{A}_{\alpha}$-vector bundles. We may endow
$\V_{\rho}\otimes\Hom_{\widetilde{A}_{\alpha}}(\V_{\rho},E)$
with structure of a~$\widetilde{G}_{\alpha}$-vector bundle. That the map
$\gamma$ is an isomorphism of vector bundles and its
$\widetilde{G}_{\alpha}$-equivariance follows from next equations.
For $(g,z)\in \widetilde{G}_{\alpha}$ we have
\begin{align*}
\gamma((g,z)\cdot(v\otimes
f))&=\gamma\big(M_{\pi(g)}\tilde{\rho}\big(\tilde{\sigma}(\pi(g))^{-1}(g,z)\big)v\big)\otimes
(\pi(g),1)\cdot f\\
&=(\pi(g),1)\cdot f
\big(M_{\pi(g)}\tilde{\rho}\big(\tilde{\sigma}(\pi(g))^{-1}(g,z)\big)v\big)\\
&=\pi(g)\bullet\!
f\big(M_{\pi(g)}\tilde{\rho}\big(\tilde{\sigma}(\pi(g))^{-1}(g,z)\big)v\big)
\!=\tilde{\sigma}(\pi(g))\!\cdot
f\big(\tilde{\rho}\big(\tilde{\sigma}(\pi(g))^{-1}(g,z)\big)v\big)\\
	&=(g,z)\cdot f(v)=(g,z)\cdot\gamma(v\otimes f).
\end{align*}
The previous argument shows that $\gamma\colon
\V_{\rho}\otimes
\Hom_{\widetilde{A}_{\alpha}}(\V_{\rho},E)\rightarrow
E$ is an isomorphism of $\widetilde{G}_{\alpha}$-vector bundles.
	
Now, if $p\colon F\rightarrow X$ is a
$\widetilde{Q}_{\beta_{\rho,\alpha}}$-equivariant vector bundle over $X$
for which the central $\S^{1}$ acts by multiplication of
scalars, then by equation~(\ref{ecu:ac2}) we know that
$\V_{\rho}\otimes F$ is a
$\big(\widetilde{G}_{\alpha},\rho\big)$-equivariant vector bundle.
The canonical isomorphism of vector bundles
\begin{align*}
F&\rightarrow \Hom_{\widetilde{A}_{\alpha}}(
\V_{\rho}, \V_{\rho}\otimes F),
\\
x&\mapsto f_{x}\colon\ v\mapsto v\otimes x
\end{align*}
is in fact $\widetilde{Q}_{\beta_{\rho,\alpha}}$-equivariant, and therefore the vector
bundles $F$ and
$\Hom_{\widetilde{A}_{\alpha}}(\V_{\rho},\V_{\rho}\otimes
F)$ are isomorphic as
$\widetilde{Q}_{\beta_{\rho,\alpha}}$-equivariant vector bundles.
	
We conclude that the inverse map of the assignment
$[E]\mapsto[\Hom_{\widetilde{A}_{\alpha}}(\V_{\rho},E)]$
is precisely the map defined by the assignment $[F]\mapsto
[\V_{\rho}\otimes F]$.
\end{proof}

%\medskip

As an immediate corollary of Theorem~\ref{teo,bij} we obtain the following
identification of the $(\widetilde{G}_{\alpha}, \rho)$-equivariant $K$-theory groups of
Definition~\ref{def,grhoalfa} with the $\beta_{\rho,\alpha}$-twisted $Q$-equivariant
$K$-theory groups provided in Definition~\ref{usualdeftwisted}.

\begin{Cor}\label{identification}
Let $X$ be a $G$-space such that $A$ acts trivially on $X$. Assume that
$\rho$ is an~$\alpha$-re\-p\-resentation of $A$ such that
$g\cdot \rho\cong \rho$ for every $g\in G$. Then the assignment
\begin{align*}
K^{\ast}_{\widetilde{G}_{\alpha}, \rho}(X)
&\stackrel{\cong}{\rightarrow} \ ^{\beta_{\rho,\alpha}}K^{\ast}_{Q}(X),
\\
[E]&\mapsto [\Hom_{\widetilde{A}_{\alpha}}(\V_{\rho},E)]
\end{align*}
defines a natural isomorphism.
\end{Cor}

Suppose now that $\alpha$ is a normalized 2-cocycle on $G$ with values in $\S^{1}$.
Consider the action of $G$ on $\Irr^{\alpha}(A)$ constructed in Section~\ref{section cocycles}.
Given $[\tau]\in \Irr^{\alpha}(A)$ let $G_{[\tau]}=\{g\in G \mid g\cdot \tau\cong \tau\}$ denote
the isotropy subgroup of the action of $G$ at $[\tau]$. The group $G_{[\tau]}$ fits into
the short exact sequence
\begin{gather*}
1\rightarrow A\rightarrow G_{[\tau]}\stackrel{\pi}{\rightarrow} Q_{[\tau]}\rightarrow 1
\end{gather*}
and $Q_{[\tau]}=G_{[\tau]}/A$ agrees with the isotropy of the group $Q$ at $[\tau]$. Let
$\big(\widetilde{G}_{[\tau]}\big)_{\alpha}$ be the central extension corresponding to the cocycle
$\alpha$ seen as a cocycle defined on $G_{[\tau]}$. Therefore we have the following commutative
diagram of central extensions
\begin{equation*}
\begin{CD}
1 @>>> \S^{1}@>>>\big(\widetilde{G}_{[\tau]}\big)_{\alpha} @>>> G_{[\tau]} @>>> 1\\
@. @V{\text{id}}VV @VVV @VVV @. \\
1 @>>> \S^{1}@>>>\widetilde{G}_{\alpha} @>>> G @>>> 1.
\end{CD}
\end{equation*}
Assume that $X$ is a compact and Hausdorff $G$-space on which $A$ acts trivially.
As before we can extend the action of $G$ on $X$ to an action of $\widetilde{G}_{\alpha}$
on $X$ in such a way that $\S^{1}$ acts trivially. Let $p\colon E\to X$ be a
$\widetilde{G}_{\alpha}$-equivariant vector bundle on which $\S^{1}$ acts by scalar
multiplication on the fibers. As $\widetilde{A}_{\alpha}$ acts trivially on
$X$ each fiber of $E$ can be seen as an $\alpha$-representation of $A$. Using
fiberwise the canonical decomposition theorem for $\alpha$-representations
(Theorem~\ref{can,decom}) we see that the assignment
\begin{align*}
\gamma\colon\ \bigoplus_{[\tau]\in \Irr^{\alpha}(A)}\V_{\tau}\otimes
\Hom_{\widetilde{A}_{\alpha}}(\V_{\tau},E)&\rightarrow E,
\\
v\otimes f&\mapsto f(v)
\end{align*}
defines an isomorphism of $\widetilde{A}_{\alpha}$-equivariant vector bundles.
Using this decomposition we obtain the following theorem.

%
%*************************************** DECOMPOSITION THEOREM ***********************************
%
\begin{Th}\label{ThD}
Under the above assumptions there is a natural isomorphism
\begin{align*}
\Psi_{X}\colon\ \phantom{i}^{\alpha}K_{G}^{\ast}(X)&\rightarrow
\bigoplus_{[\tau]\in G\backslash \Irr^{\alpha}(A)}
K^{\ast}_{(\widetilde{G}_{[\tau]})_{\alpha},\tau}(X),
\\
[E]&\mapsto \bigoplus_{[\tau]\in G\backslash \Irr^{\alpha}(A)}
\big[\V_{\tau}\otimes\Hom_{\widetilde{A}_{\alpha}}(\V_{\tau},E)\big].
\end{align*}
This isomorphism is
functorial on maps $X\rightarrow Y$ of $G$-spaces on which $A$ acts trivially.
\end{Th}

\begin{proof}
The proof of this theorem follows the same lines of the proof of~\cite[Theorem 3.1]{GomezUribe}
so we only provide an outline of the proof.

Let us show first that the map $\Psi_{X}$ is well defined. To see this we have to
show that if $\rho$ is an $\alpha$-representation of $A$ then
$\V_{\rho}\otimes\Hom_{\widetilde{A}_{\alpha}}(\V_{\rho},E)$
has the structure of a $\big(\big(\widetilde{G}_{[\rho]}\big)_{\alpha},\rho\big)$-vector bundle.
Following the notation of Section~\ref{section cocycles} fix a set theoretical
section $\sigma\colon Q_{[\rho]}\to G_{[\rho]}$ for the projection map $\pi\colon G_{[\rho]}\to Q_{[\rho]}$
in such a way that $\sigma(1)=1$. Also, for every $q\in Q_{[\rho]}$ fix an~element
$M_{q}\in U(V_{\rho})$ such that
\begin{gather*}
\sigma(q)\cdot \rho(a) =\alpha\big(\sigma(q)^{-1}a,\sigma(q)\big)\alpha\big(\sigma(q),\sigma(q)^{-1}a\big)^{-1}
\rho\big(\sigma(q)^{-1}a\sigma(q)\big)=M_{q}^{-1}\rho(a)M_{q}.
\end{gather*}
This is possible as $\sigma(q)\in G_{[\rho]}$ so that we have $\sigma(q)\cdot \rho\cong \rho$.
We can choose $M_{1}=1$ as $\sigma(1)=1$.

Now if $(h,z)\in \big(\widetilde{G}_{[\rho]}\big)_{\alpha}$ and $v\otimes f\in
\V_{\rho}\otimes \Hom_{\widetilde{A}_{\alpha}}(\V_{\rho},E)$
we define $M_{(h,z)}\in U(V_{\rho})$ by
\begin{gather*}
M_{(h,z)}:=M_{\pi(h)}\tilde{\rho}\big(\tilde{\sigma}(\pi(h))^{-1}(h,z)\big),
\end{gather*}
where $\tilde{\rho}$ and $\tilde{\sigma}$ are defined in a similar way as in
Theorem~\ref{teo,bij}. Observe that $M_{(a,z)}=\tilde{\rho}(a,z)$ for all
$(a,z)\in \widetilde{A}_{\alpha}$. Moreover, given $(h,z)\in \big(\widetilde{G}_{[\rho]}\big)_{\alpha}$
and $v\otimes f\in \V_{\rho}\otimes \Hom_{\widetilde{A}_{\alpha}}(\V_{\rho},E)$
we define
\begin{gather*}
(h,z)\star(v\otimes f)=M_{(h,z)}v\otimes (h,z)\bullet f,
\end{gather*}
where $(h,z)\bullet f (w)=(h,z)f\big(M_{(h,z)}^{-1}w\big)$. Unraveling the definitions it
can be seen that this defines an action of
$(\widetilde{G}_{[\rho]})_{\alpha}$ on $\V_{\rho}\otimes \Hom_{\widetilde{A}_{\alpha}}(\V_{\rho},E)$
in such a way that the central factor $\S^{1}$ acts by multiplication of scalars.
This way $\V_{\rho}\otimes \Hom_{\widetilde{A}_{\alpha}}(\V_{\rho},E)$ has the structure
of a $\big(\widetilde{G}_{[\rho]}\big)_{\alpha}$-vector bundle and $A$ acts by the
$\alpha$-representation $\rho$ on the fibers so that
$\big[\V_{\rho}\otimes \Hom_{\widetilde{A}_{\alpha}}(\V_{\rho},E)\big]\in
K^{\ast}_{(\widetilde{G}_{[\rho]})_{\alpha},\rho}(X)$. This shows that
$\Psi_{X}$ is well defined.


Next we show that $\Psi_{X}$ is an isomorphism. For this write $\Irr^{\alpha}(A)=
\mathcal{A}_{1}\sqcup\mathcal{A}_{2}\sqcup\dots \sqcup\mathcal{A}_{k}$, where
$\mathcal{A}_{1},\mathcal{A}_{2},\dots, \mathcal{A}_{k}$ are the different
$G$-orbits of the action of $G$ on $\Irr^{\alpha}(A)$ defined in equation~(\ref{def:act}).
For every $1\leq i\leq k$ define
\begin{gather*}
E_{\mathcal{A}_{i}}=\bigoplus_{[\tau]\in \mathcal{A}_{i}}
\V_{\tau}\otimes \Hom_{\widetilde{A}_{\alpha}}(\V_{\tau},E).
\end{gather*}
Note that $\V_{\tau}\otimes \Hom_{\widetilde{A}_{\alpha}}(\V_{\tau},E)$ is an
$\widetilde{A}_{\alpha}$-equivariant vector bundle over $X$, so each
$E_{\mathcal{A}_{i}}$ is also an $\widetilde{A}_{\alpha}$-equivariant vector bundle over $X$
and the map
\begin{gather*}
\gamma\colon\ \bigoplus_{i=1}^{k}E_{\mathcal{A}_{i}}=\bigoplus_{[\tau]\in \Irr^{\alpha}(A)}
\V_{\tau}\otimes \Hom_{\widetilde{A}_{\alpha}}(\V_{\tau},E)\rightarrow E
\end{gather*}
defines an isomorphism of $\widetilde{A}_{\alpha}$-vector bundles. We are going to show that each
$E_{\mathcal{A}_{i}}$ is a $\widetilde{G}_{\alpha}$-vector bundle and that the map $\gamma$ is
$\widetilde{G}_{\alpha}$-equivariant. For this fix an index $1\leq i\leq k$ and an irreducible
$\alpha$-representation $\rho\colon A\rightarrow U(V_{\rho})$ such that
$[\rho]\in\mathcal{A}_{i}$. The elements in $\mathcal{A}_{i}$ can be written in the
form $[g_{1}\cdot\rho], \dots, [g_{r_{i}}\cdot\rho]$ for some elements
$g_{1}=1,g_{2},\dots,g_{r_{i}}\in G$. Therefore
\begin{gather*}
E_{\mathcal{A}_{i}}=\bigoplus_{j=1}^{r_{i}}\V_{g_{j}\cdot\rho}\otimes
\Hom_{\widetilde{A}_{\alpha}}(\V_{g_{j}\cdot\rho},E).
\end{gather*}
We can give a structure of
$\widetilde{G}_{\alpha}$-space on $E_{\mathcal{A}_{i}}$ in the following way.
Suppose that $(g,s)\in\widetilde{G}_{\alpha}$ and~that
$v\otimes f\in \V_{\rho}\otimes \Hom_{\widetilde{A}_{\alpha}}(\V_{\rho},E)_{x}$.
Decompose $gg_{j}$ in the form $gg_{j}=g_{l}h$, where $1\leq l\leq r_{i}$ and $h\in G_{[\rho]}$.
In other words $g_{l}$ is the representative chosen for the coset
$(gg_{j})G_{[\rho]}$ and $h=g_{l}^{-1}gg_{j}$. Then
\begin{gather*}
(g,s)(g_{j},1)=(g_{l},1)(h,z),
\end{gather*}
where $z=s\alpha(g,g_{j})\alpha(g_{l},h)^{-1}$. We define
\begin{gather*}
(g,s)\star (v\otimes f):=M_{(h,z)}v\otimes (g,s)\bullet
f \in \big(\V_{g_{l}\cdot\rho}\otimes
\Hom_{\widetilde{A}_{\alpha}}(\V_{g_{l}\cdot\rho},E)\big)_{(g,s)\cdot x},
\end{gather*}
where
\begin{gather*}
(g,s)\bullet f(w)=(g,s)f\big(M_{(h,z)}^{-1}w\big)
=(g,s)(h,z)^{-1}\tilde{\sigma}(\pi(h))f\big(M_{\pi(h)}^{-1}w\big).
\end{gather*}
It can be seen that this defines an action of $\widetilde{G}_{\alpha}\in E_{\mathcal{A}_{i}}$
for each $1\leq i\leq k$ making the vector bundle $E_{\mathcal{A}_{i}}$
into a $\widetilde{G}_{\alpha}$-vector bundle in such a way that the central factor
$\S^{1}$ acts by multiplication of scalars. Furthermore, the map
\begin{gather*}
\gamma\colon\ \bigoplus_{i=1}^{k}E_{\mathcal{A}_{i}} \rightarrow E
\end{gather*}
is an isomorphism of $\widetilde{A}_{\alpha}$-vector and the map $\gamma$ is
$\widetilde{G}_{\alpha}$-equivariant so that $\gamma$ is an isomorphism of~$\widetilde{G}_{\alpha}$-vector bundles. Now, the desired map $\Psi_{X}$ can be seen as the
direct sum $\oplus_{i=1}^{k}\Psi_{X}^{i}$ choosing for each $\mathcal{A}_{i}$ a representation
$[\tau_{i}]\in\mathcal{A}_{i}$, where each $\Psi_{X}^{i}$ is given by
\begin{align*}
\Psi_{X}^{i}\colon \phantom{i}^{\alpha}K_{G}^{\ast}(X)&\rightarrow
K^{*}_{\left(\widetilde{G}_{[\tau_{i}]}\right)_{\alpha},\tau_{i}}(X),
\\
[E]&\mapsto \big[\V_{\tau_{i}}\otimes\Hom_{\widetilde{A}_{\alpha}}(\V_{\tau_{i}},E)\big].
\end{align*}
In what follows we will construct the map
$\zeta_{i}\colon K^{*}_{\left(\widetilde{G}_{[\tau_{i}]}\right)_{\alpha},\tau_{i}}(X)
\rightarrow \phantom{i}^{\alpha}K_{G}^{\ast}(X)$
which will be the right inverse of $\Psi_{X}^{i}$. Take $\rho=\tau_{i}$ and consider
a vector bundle $F\in\text{Vect}_{\left(\widetilde{G}_{[\rho]}\right)_{\alpha},\rho}(X)$.
Fix $g_{1}=1,g_{2},\dots,g_{r}$ representatives for the different cosets in $G/G_{[\rho]}$.
Let
\begin{gather*}
L_{F}:=\bigoplus_{j=1}^{r}\big[(g_{j},1)^{-1}\big]^{\ast} F.
\end{gather*}
Using ideas similar to the ones used in~\cite[Theorem 3.1]{GomezUribe} we can endow $L_{F}$
with the structure of a $\widetilde{G}_{\alpha}$-vector bundle on which the
central factor $\S^{1}$ acts by multiplication of scalars in such a~way that the action
of $\big(\widetilde{G}_{[\rho]}\big)_{\alpha}$ on $\big[(g_{1},1)^{-1}\big]^{\ast} F\cong F$ agrees with
the given action of $\big(\widetilde{G}_{[\rho]}\big)_{\alpha}$ on~$F$. We~define
\begin{align*}
\zeta_{i}\colon\ K^{*}_{\left(\widetilde{G}_{[\tau_{i}]}\right)_{\alpha},\tau_{i}}(X)
&\rightarrow \phantom{i}^{\alpha}K_{G}^{\ast}(X),
\\
[F]&\mapsto [L_{F}].
\end{align*}
Now, since we have the isomorphism
$\V_{\rho}\otimes \Hom_{\widetilde{A}_{\alpha}}
\big(\V_{\rho},\oplus_{j=1}^{r}\big[(g_{j},1)^{-1}\big]^{\ast} F\big)\cong F$ as
$\big(\widetilde{G}_[\rho]\big)_{\alpha}$ vector bundles. We obtain at the level of $K$-theory that
$\zeta_{i}$ is a right inverse for $\Psi_{i}$ so that
$\zeta=\oplus_{i=1}^{r}\zeta_{i}$ is a right inverse for $\Psi_{X}$. In a similar way, using the
work given above it can be seen that $\zeta$ is also a left inverse so that the
map $\Psi_{X}$ is indeed an isomorphism.

To finish, we observe that functoriality follows from the fact that if $\tau$ is an
$\alpha$-representation of $A$ then the bundles
$\V_{\tau}\otimes \Hom_{\widetilde{A}_{\alpha}}(\V_{\tau},f^{\ast} E)$ and
$f^{\ast}\big(\V_{\tau}\otimes \Hom_{\widetilde{A}_{\alpha}}(\V_{\tau},E)\big)$
are canonically isomorphic as $((\widetilde{G}_{[\tau]})_{\alpha},\tau)$-equivariant bundles
whenever $f\colon Y\rightarrow X$ is a
$G$-equivariant map from spaces on which $A$ acts trivially.
\end{proof}

%\medskip

As a result of Theorem~\ref{ThD} and Corollary~\ref{identification} we obtain the following
theorem that is the main result of this article.

\begin{Th}\label{Thmain}
Suppose that $A$ is a normal subgroup of a finite group $G$. Let $\alpha$
be a normalized 2-cocycle on $G$ with values in $\S^{1}$ and $X$ a
compact $G$-space on which $A$ acts trivially.
Then there is a natural isomorphism
\begin{align*}
\Phi_{X}\colon\ \phantom{i}^{\alpha}K_{G}^{\ast}(X)&\rightarrow
\bigoplus_{[\tau]\in G\backslash \Irr^{\alpha}(A)}\phantom{i}
^{\beta_{\tau,\alpha}}K^{\ast}_{Q_{[\tau]}}(X),
\\
[E]&\mapsto \bigoplus_{[\tau]\in G\backslash \Irr^{\alpha}(A)}
\big[\Hom_{\widetilde{A}_{\alpha}}(\V_{\tau},E)\big].
\end{align*}
Here $\beta_{\tau,\alpha}$ is the $2$-cocycle associated to $\tau$ and $\alpha$ as defined in equation~$(\ref{defQbeta})$. This isomorphism is functorial on maps $X\rightarrow Y$ of
$G$-spaces on which $A$ acts trivially.
\end{Th}

%
%*************************************** SECTION AHSS ***********************************
